\documentclass[11pt,a4paper]{article}

\usepackage[utf8]{inputenc}
\usepackage[T1]{fontenc}
\usepackage[polish,english]{babel}
\usepackage{geometry}
\usepackage{amsmath,amssymb,amsthm,mathtools}
\usepackage{bm}
\usepackage{graphicx}
\usepackage{booktabs}
\usepackage{xcolor}
\usepackage{hyperref}
\usepackage{listings}
\usepackage{enumitem}
\usepackage{microtype}
\usepackage{fancyhdr}

\geometry{margin=25mm}
\hypersetup{colorlinks=true, linkcolor=blue, urlcolor=blue, citecolor=blue}
\setlength{\headheight}{14pt}
\pagestyle{fancy}
\fancyhf{}
\lhead{Hamiltonian fazowo-intencjonalny}
\rhead{podkladka formalno-techniczna}
\cfoot{\thepage}

\lstset{
  basicstyle=\ttfamily\small,
  breaklines=true,
  frame=single,
  columns=fullflexible,
  keepspaces=true,
  showstringspaces=false,
  tabsize=2,
  language=Python
}

\newcommand{\Iop}{\hat{\mathcal I}}
\newcommand{\Wop}{\hat W}
\newcommand{\Abe}{\mathcal A^{ABE}}
\newcommand{\Fbe}{\mathcal F^{BE}}
\newcommand{\Q}{\mathcal Q}
\newcommand{\Hs}{\hat H_s}
\newcommand{\kappaI}{\kappa_I}
\newcommand{\dd}{\,\mathrm d}

\title{\textbf{Hamiltonian fazowo-intencjonalny}\\[2mm]
\large Aharonov--Bohm, Berry, Euler, Intention Operator, rytm Collatza i seed free twin primes}
\author{Adrian Lipa / USER007 \\ opracowanie robocze: Zosia / Codex GPT}
\date{19 czerwca 2026}

\begin{document}
\selectlanguage{polish}
\maketitle

\begin{abstract}
Ta nota zbiera w jedna, czytelna podkladke robocza model, w ktorym czesc swobodna Hamiltonianu nie jest klasycznym $H_0$ czastki swobodnej, lecz fazowym generatorem intencji/informacji. Geometria fazy jest modelowana jako polaczenie $U(1)$ z wkladami Aharonova--Bohma, Berry'ego i Eulera. Generator jest dany przez Operator Intencji/Informacji, skalowany stala $\ln 2/(24\pi)$, taktowany dyskretnym rytmem Collatza/Kolaca i inicjowany seedem typu free twin primes. Dodatkowo wprowadzamy Euler--Berry constraint, ktory wiaze defekt fazowy $\epsilon$, sektor $D$ oraz domkniecie holonomii. Dokument zawiera rowniez minimalne symulacje referencyjne i kod do dalszej implementacji.
\end{abstract}

\section*{Status i granice}

\begin{itemize}[leftmargin=*]
\item \textbf{Status formalny:} ansatz matematyczny, scaffold symulacyjny, material do dalszej implementacji.
\item \textbf{Nie jest to:} eksperymentalny dowod fizyczny, zewnetrzny NOEMA CURRENT/ACK ani kanoniczny wynik bez osobnego receiptu.
\item \textbf{Zasada robocza:} faza jest traktowana jako pierwotny kanal dynamiki, a nie jako mala poprawka do czastki swobodnej.
\item \textbf{Granica notacji:} symbol $D$ jest tutaj uzyty jako sektor Eulera/Berry'ego: stopien, indeks defektu, liczba topologiczna albo efektywny wymiar/domkniecie. Jesli w kanonie projektu $D$ ma wezszke znaczenie, nalezy podstawic tamta definicje, zachowujac rownanie constraintu.
\end{itemize}

\tableofcontents
\newpage

\section{Glowna teza}

W standardowym rozbiciu perturbacyjnym pisze sie zwykle
\begin{equation}
H = H_{\mathrm{free}} + H_{\mathrm{int}}.
\end{equation}
W modelu relacyjno-fazowym to rozbicie jest niewlasciwe. Tutaj czesc swobodna nie oznacza czastki bez oddzialywania, lecz swobodny generator fazy:
\begin{equation}
\boxed{H = H_{\Phi,\mathcal I,C,s} + H_{\perp}.}
\end{equation}
Gdzie:
\begin{itemize}[leftmargin=*]
\item $H_{\Phi,\mathcal I,C,s}$ - fazowy free part: Intention/Information Operator, rytm Collatza, seed twin-prime;
\item $H_{\perp}$ - reszta: dynamika horyzontalna, ksztalt, amplituda, potencjal, deformacje wzgledem polaczenia fazowego;
\item $s=(p,p+2)$ - seed free twin primes.
\end{itemize}

Najkrotsza postac robocza to:
\begin{equation}
\boxed{
\Hs(\tau,k)
=
\frac{\hbar}{\Delta\tau_k}\rho_s(k)
\left[\frac{\ln2}{24\pi}\Wop_s+\delta\Iop_{0}(\tau)\right]
+
\frac12 g^{ab}\hat\Pi_a^{(s,k)}\hat\Pi_b^{(s,k)}
+V_s(q,\tau,k).}
\label{eq:main_hamiltonian}
\end{equation}

\section{Przestrzen relacyjna i wiazka fazowa}

Niech $\Q$ bedzie relacyjna przestrzenia konfiguracji, a $q^a$ jej wspolrzednymi. Faza jest wspolrzedna wlokna $U(1)$:
\begin{equation}
U(1) \hookrightarrow P \rightarrow \Q.
\end{equation}
Stan ma postac sekcji:
\begin{equation}
\Psi_s(q,\chi,k) \in L^2(\Q,L_s),
\end{equation}
gdzie $L_s$ jest wiazka liniowa zależna od seeda twin-prime $s$.

Wprowadza sie kowariantna pochodna fazy:
\begin{equation}
D_t\chi = \dot\chi + \Abe_a(q)\dot q^a.
\end{equation}
Faza nie jest wiec izolowana zmienna zegarowa, lecz jest mierzona wzgledem polaczenia holonomicznego.

\section{Polaczenie Aharonov--Bohm--Berry--Euler}

Definiujemy laczne polaczenie:
\begin{equation}
\boxed{
\Abe
=
\mathcal A_{AB}+\mathcal A_B+\mathcal A_E.}
\end{equation}
Skladniki:
\begin{align}
\mathcal A_{AB} &= \frac{e}{\hbar}A_i(x)\dd x^i,\\
\mathcal A_B &= i\langle u(q)|\dd u(q)\rangle,\\
\mathcal A_E &= s_E\omega_E.
\end{align}
Dla sektora spinowego $1/2$ przyjmujemy:
\begin{equation}
s_E=\frac12.
\end{equation}

Holonomia po petli $\gamma$:
\begin{equation}
U_\gamma
=
\exp\left(i\oint_\gamma \Abe\right).
\end{equation}

Calkowita faza geometryczna:
\begin{equation}
\Phi_{ABE}[\gamma]
=
\Phi_{AB}+\Phi_B+\Phi_E
=
\oint_\gamma \Abe.
\end{equation}

\section{Operator Informacji i Operator Intencji}

Stala informacyjna projektu:
\begin{equation}
\boxed{\kappaI = \frac{\ln2}{24\pi}.}
\end{equation}
Sama stala nie jest jeszcze pelnym operatorem. Jako operator skaluje strukture windingu/intencji:
\begin{equation}
\boxed{\Iop_s = \kappaI \Wop_s.}
\end{equation}
Po dodaniu dynamiki z zera:
\begin{equation}
\boxed{\Iop_s(\tau)=\kappaI\Wop_s+\delta\Iop_0(\tau).}
\end{equation}

Czesc hermitowska, jesli pelny operator ma komponent selekcyjny lub dyssypatywny:
\begin{equation}
\Iop_{H} = \frac12\left(\Iop+\Iop^\dagger\right).
\end{equation}
Hamiltonian jako generator unitarnej fazy powinien uzywac czesci samosprzezonej.

\subsection{Dziedzina operatora}

Dziedzina operatora to nie wartosc stalej $\ln2/(24\pi)$, lecz zbior stanow, na ktorych operator jest legalnie okreslony:
\begin{equation}
D(\Iop_s)=
\left\{
\Psi_s\in L^2(\Q,L_s):
\Wop_s\Psi_s\in L^2(\Q,L_s),\quad
\Psi_s\text{ spelnia warunki holonomii}
\right\}.
\end{equation}
Warunek holonomii:
\begin{equation}
\Psi_s(\gamma(1))
=
\exp\left(i\Phi_{ABE}[\gamma]\right)
\Psi_s(\gamma(0)).
\end{equation}

\section{Twin-prime seed i rytm Collatza/Kolaca}

Seed jest para blizniaczych liczb pierwszych:
\begin{equation}
\boxed{s=(p,p+2),\qquad p,p+2\in\mathbb P.}
\end{equation}
Mapa Collatza:
\begin{equation}
C(n)=
\begin{cases}
n/2, & n\equiv 0\pmod 2,\\
3n+1, & n\equiv 1\pmod 2.
\end{cases}
\end{equation}
Orbita pary:
\begin{equation}
\mathcal O_s(k)=\left(C^k(p),C^k(p+2)\right).
\end{equation}

Rytm Collatza jest modularem modelu. W symulacji referencyjnej uzywamy ograniczonego rytmu logarytmicznego:
\begin{equation}
\rho_s(k)=1+\alpha\tanh\left[\frac12\left(
\ln\frac{C^{k+1}(p)}{C^k(p)}+
\ln\frac{C^{k+1}(p+2)}{C^k(p+2)}
\right)\right].
\end{equation}
To jest wybor testowy, nie jedyna mozliwa definicja.

\section{Fazowy free part Hamiltonianu}

Jeden krok fazowy:
\begin{equation}
U_{\Phi}(k+1,k)
=
\exp\left[-i\rho_s(k)\Iop_s(k)\right].
\end{equation}
Po wprowadzeniu metaczasu $\Delta\tau_k$:
\begin{equation}
U_{\Phi}(k+1,k)
=
\exp\left[-\frac{i}{\hbar}\Delta\tau_k H_{\Phi,\mathcal I,C,s}(k)\right].
\end{equation}
Zatem:
\begin{equation}
\boxed{
H_{\Phi,\mathcal I,C,s}(\tau,k)
=
\frac{\hbar}{\Delta\tau_k}\rho_s(k)
\left[\kappaI\Wop_s+\delta\Iop_0(\tau)\right].}
\end{equation}
To jest wlasciwy sens zdania:
\begin{center}
\emph{free part Hamiltonianu jest faza generowana intencja/informacja, nie czastka swobodna.}
\end{center}

\section{Pedy kowariantne i czesc horyzontalna}

Kowariantny ped relacyjny:
\begin{equation}
\boxed{
\hat\Pi_a^{(s,k)}
=
-i\hbar\nabla_a
-\hbar\Abe_a(q)
-\lambda_s(k)\mathcal I_a(q).}
\end{equation}
Wtedy:
\begin{equation}
H_{\perp,s}(k)
=
\frac12g^{ab}\hat\Pi_a^{(s,k)}\hat\Pi_b^{(s,k)}+V_s(q,k).
\end{equation}

Pelny Hamiltonian jest rownaniem \eqref{eq:main_hamiltonian}. Jego interpretacja:
\begin{equation}
\boxed{
\text{seed}\rightarrow\text{Collatz rhythm}\rightarrow\text{Intention/Information generator}\rightarrow\text{phase}\rightarrow\text{holonomic relational motion}.}
\end{equation}

\section{Euler--Berry constraint: zwiazanie $\epsilon$, $D$ i fazy}

Tu wprowadzamy brakujacy wezly constraintu. Niech $\Sigma$ bedzie powierzchnia z brzegiem $\partial\Sigma=\gamma$. Krzywizna Berry--Euler:
\begin{equation}
\Fbe = \dd\mathcal A_B+s_E\dd\omega_E.
\end{equation}
Definiujemy sektor:
\begin{equation}
\boxed{
D_\Sigma
=
\frac{1}{2\pi}\int_\Sigma \Fbe.}
\end{equation}
Gdy $D_\Sigma$ jest topologicznie skwantowany, dziala jak indeks defektu/stopien/domkniecie. Faza Euler--Berry jest wtedy:
\begin{equation}
\Phi_{BE}=2\pi D_\Sigma.
\end{equation}
Jesli wydzielamy Euler z $s_E=1/2$, mozna rowniez pisac:
\begin{equation}
\Phi_E=2\pi s_E D_E.
\end{equation}

Calkowita faza modelu:
\begin{equation}
\Phi_{\mathrm{tot}}
=
\Phi_{AB}+2\pi D_\Sigma+\Phi_{\mathcal I,C,s},
\end{equation}
gdzie:
\begin{equation}
\Phi_{\mathcal I,C,s}
=
\sum_{k=0}^{K-1}\rho_s(k)\langle\Iop_s(k)\rangle.
\end{equation}

Defekt domkniecia fazy:
\begin{equation}
\boxed{
\epsilon_{\gamma,\Sigma}
=
\operatorname{dist}_{\mathbb Z}\left(
\frac{1}{2\pi}
\left[
\Phi_{AB}+2\pi D_\Sigma+\Phi_{\mathcal I,C,s}
\right]
\right).}
\end{equation}
Tu $\operatorname{dist}_{\mathbb Z}(x)$ oznacza odleglosc $x$ od najblizszej liczby calkowitej, ze znakiem albo bez znaku, zależnie od wariantu.

Euler--Berry constraint:
\begin{equation}
\boxed{
\epsilon_{\gamma,\Sigma}=0
\quad\text{albo roboczo}\quad
|\epsilon_{\gamma,\Sigma}|\le\epsilon_\star.}
\end{equation}

To jest formalny zapis zdania: \textbf{Euler--Berry constraint wiaze epsilon, $D$ i ogolna faze}. Jesli $D$ zmienia sie o jeden sektor, faza przesuwa sie o kwant $2\pi$ albo $\pi$ w sektorze spinowym $1/2$, a dopuszczalnosc stanu zalezy od tego, czy calosc domyka holonomie.

\section{Lagrangian relacyjno-fazowy}

Odpowiadajacy Lagrangian moze miec postac:
\begin{equation}
L
=
\frac12g_{ab}(q)\dot q^a\dot q^b
+
\frac{I_\phi}{2}\left(\dot\chi+\Abe_a\dot q^a\right)^2
+J_0\left(\dot\chi+\Abe_a\dot q^a\right)
-V(q)
+L_{\mathcal I,C,s}.
\end{equation}
Skladnik intencyjny:
\begin{equation}
L_{\mathcal I,C,s}
\sim
\hbar\rho_s(k)\langle\Iop_s(k)\rangle.
\end{equation}
Transformata Legendre'a daje fazowy generator Hamiltonianu, jesli czesc fazowa ma dobrze zdefiniowany moment sprzezony albo jest ujęta jako Floquetowski generator krokowy.

\section{Symulacje referencyjne}

Symulacje nie sa dowodem fizycznym. Ich cel jest techniczny:
\begin{enumerate}[leftmargin=*]
\item sprawdzic, czy rytm Collatza z twin-prime seeda daje dobrze policzalna faze;
\item przeskanowac Euler--Berry closure error $\epsilon(D)$;
\item sprawdzic hermitowskosc i unitarna ewolucje w skonczonym modelu macierzowym;
\item stworzyc punkt startowy dla implementacji w NOEMA/ApocalypOS bez mieszania tego z kanonicznym SoT.
\end{enumerate}

\subsection{Wyniki uruchomionej paczki}

Konfiguracja testowa:
\begin{itemize}[leftmargin=*]
\item seed: $s=(11,13)$;
\item liczba krokow Collatza: $80$;
\item $\kappaI=\ln2/(24\pi)\approx 0.009193150006360484$;
\item scan $D\in[-8,8]$;
\item macierzowy test Hamiltonianu: wymiar $8$, liczba krokow $60$.
\end{itemize}

Wybrane wyniki:
\begin{table}[h]
\centering
\begin{tabular}{ll}
\toprule
Wielkosc & Wartosc \\
\midrule
Najlepszy sektor $D$ w skanie & $-5$ \\
Najmniejszy $|\epsilon|$ w skanie & $0.2135735933580536$ \\
Koncowa faza intencyjna & $-0.09192134655845716$ \\
Maksymalny blad hermitowskosci & $0.0$ \\
Maksymalne odchylenie normy od $1$ & $2.886579864025407\times10^{-15}$ \\
\bottomrule
\end{tabular}
\caption{Wyniki referencyjne wygenerowane przez dolaczone skrypty.}
\end{table}

Interpretacja: macierzowy scaffold zachowuje unitarna norme numerycznie; skan $D$ pokazuje, ze przy arbitralnych fazach testowych $\Phi_{AB}=0.73$, $\Phi_B=1.17$ i przy prostej definicji rytmu nie uzyskano idealnego domkniecia $\epsilon=0$, ale uzyskano mierzalny blad domkniecia, ktory mozna minimalizowac przez wybor sektora $D$ albo przez kalibracje faz.

\subsection{Wizualizacja fazy intencyjnej}

\begin{figure}[h]
\centering
\includegraphics[width=0.82\textwidth]{figures/collatz_intention_phase.png}
\caption{Skumulowana faza intencyjna z twin-prime seeda $(11,13)$ i rytmu Collatza.}
\end{figure}

\begin{figure}[h]
\centering
\includegraphics[width=0.82\textwidth]{figures/euler_berry_epsilon_scan.png}
\caption{Skan defektu domkniecia $|\epsilon|$ wzgledem sektora Euler--Berry $D$.}
\end{figure}

\begin{figure}[h]
\centering
\includegraphics[width=0.82\textwidth]{figures/unitary_norm_check.png}
\caption{Kontrola normy w modelu macierzowym: ewolucja przez hermitowski Hamiltonian fazowo-horyzontalny.}
\end{figure}

\begin{figure}[h]
\centering
\includegraphics[width=0.82\textwidth]{figures/energy_expectation_demo.png}
\caption{Przykladowa wartosc oczekiwana energii przy czasowo zaleznym generatorze fazy.}
\end{figure}

\clearpage
\section{Minimalny program dalszej pracy}

\begin{enumerate}[leftmargin=*]
\item Ustalic kanoniczna definicje $\Wop_s$: winding, signed Collatz imbalance, operator orbitalny, albo operator na grafie NOEMA.
\item Ustalic kanoniczna definicje rytmu $\rho_s(k)$: parzystosc, log-kontrakcja, czas stopu, lokalny gradient, albo mieszany regulator.
\item Ustalic, czym dokladnie w kanonie jest $D$: indeks Eulera, wymiar, defect charge, degree, divisor, sektor spinowy, czy efektywny operator.
\item Przepisac Euler--Berry constraint jako test Doctor/Guard: $|\epsilon|\le\epsilon_\star$.
\item Zrobic wersje HDF5/NOEMA Surface: wyniki symulacji jako append-only packet, nie jako fake CURRENT.
\item Dodac test samosprzezonosci $\Iop_H$ i test zachowania normy dla kazdej przyszlej wersji Hamiltonianu.
\end{enumerate}

\appendix
\clearpage
\section{Kod: Collatz phase i Euler--Berry epsilon scan}

\lstinputlisting{scripts/collatz_phase_sim.py}

\clearpage
\section{Kod: macierzowy Hamiltonian i test unitarnosci}

\lstinputlisting{scripts/hamiltonian_matrix_demo.py}

\clearpage
\section{Kod: runner symulacji}

\lstinputlisting{scripts/run_all.py}

\clearpage
\section*{NOEMA writeback plan}

\begin{verbatim}
NOEMA_WRITEBACK_PLAN:
write_required: true
new_snap_type: formal_technical_podkladka
append_only: true
update_current_pointer: false
artifacts_created:
  - phase_hamiltonian_podkladka.pdf
  - phase_hamiltonian_podkladka_source.zip
blockers:
  - brak dostepu do realnego NOEMA SoT/write executor
  - dokument jest artifact-grounded, nie external CURRENT/ACK
\end{verbatim}

\end{document}
